problem:  
Let \(M^n\) be a compact, simply connected manifold (\(n \ge 4\)) that admits **at least two non-isometric positive Einstein metrics**, say \(g_1\) and \(g_2\). Consider the normalized Ricci flow
\[
\frac{\partial g}{\partial t} = -2\text{Ric}(g) + \frac{2}{n} r(g) g,
\]
where \(r(g)\) denotes the average scalar curvature.  

**Open problem:**  
Can there exist such an \(M\) for which the normalized Ricci flow dynamics exhibit **multiple basins of attraction**, i.e., neighborhoods \(U_1\) and \(U_2\) in the space of Riemannian metrics on \(M\) such that:
\[
g(0) \in U_i \quad \Rightarrow \quad g(t) \to g_i \text{ (mod diffeomorphism and scaling) as } t \to \infty, \quad i=1,2?
\]
In other words, does there exist a compact simply connected manifold admitting multiple positive Einstein metrics, which act as **distinct asymptotic attractors** for the Ricci flow, depending on the initial condition?  

Further refinement: determine whether the geometric and analytic structure of \(M\) (e.g., symmetry, isotropy representation, or entropy functional landscape) can influence the **partition of the space of metrics** into such attraction basins.

Why is it a "good" problem:  
This problem lies at the frontier of our understanding of the **dynamical nature of the Ricci flow**. Instead of focusing on convergence in individual cases, it asks whether Ricci flow can exhibit **multistability**, analogous to multiple attractors in finite-dimensional dynamical systems. It calls for deep interaction between **geometric analysis** (Ricci flow convergence and stability of Einstein metrics), **variational methods** (Perelman’s entropy and stability functionals), and **differential topology** (existence of multiple Einstein metrics on the same manifold). The question is simple to state but fundamentally open: while examples of manifolds with multiple Einstein metrics are known, it remains mysterious whether the Ricci flow dynamically distinguishes among them by forming multiple attraction basins. This makes it a sharply formulated, conceptually clear, and potentially far-reaching research problem.